\documentclass[a4paper]{article}

\IfFileExists{INTERSPEECH2021.sty}{%
  \usepackage{INTERSPEECH2021}%
}{%
  \IfFileExists{report_template/INTERSPEECH2021.sty}{%
    \usepackage{report_template/INTERSPEECH2021}%
  }{%
    \usepackage{../report_template/INTERSPEECH2021}%
  }%
}

\usepackage{amsmath}
\usepackage{booktabs}

\title{Deliveroo.js}
\name{Riccardo Miolato, Matteo Morellini (mat. 268427)} 

\address{University of Trento}
\email{riccardo.miolato@studenti.unitn.it, matteo.morellini@studenti.unitn.it}

\begin{document}

\maketitle

\section{Cluster-Based Search with Resumable Coverage}
\label{subsec:cluster-search}

When no parcel with positive expected utility is known, the agent searches the
parcel-spawning cells. Rather than selecting an individual cell at random,
spawning cells are grouped into clusters, and the agent plans a route that
guarantees complete visual coverage of one cluster at a time.

This strategy reflects the stochastic parcel-generation mechanism of the
environment. At each generation event, provided that the global parcel limit
has not been reached, the simulator selects uniformly at random one currently
empty spawning cell. If $\mathcal{E}_t\subseteq\mathcal{S}$ is the set of
eligible spawning cells at time $t$, then

\[
    \Pr(c\text{ is selected}\mid c\in\mathcal{E}_t)
    = \frac{1}{|\mathcal{E}_t|}.
\]

There is no additional location-dependent weighting. Nevertheless, a region
containing more eligible spawning cells has a proportionally greater aggregate
probability of receiving a parcel. Systematically revisiting and completely
observing spawning clusters is therefore preferable to repeatedly moving
towards independently sampled cells, since it provides spatial coverage while
accounting for the changing set of eligible locations.

\subsection{Cluster Selection}

Let $\mathcal{S}$ denote the set of parcel-spawning cells. Clusters
$C_1,\ldots,C_k$ are the connected components of $\mathcal{S}$ under
four-directional adjacency; diagonal adjacency alone does not connect two
cells.

For each cluster $C_i$, the agent stores its last complete-observation time
$t_i$. When no scan is active, it selects the least recently completed cluster,

\[
    C^* = \arg\min_{C_i} t_i,
\]

where unvisited clusters take priority over previously completed ones. Ties
are resolved in favour of the cluster closest to the agent.

\subsection{Visibility Coverage}

With Manhattan observation radius $R$, a spawning cell $c$ is visible from
position $p$ whenever

\[
    d_M(p,c)=|p_x-c_x|+|p_y-c_y| \leq R.
\]

A route $\pi$ covers a cluster $C_i$ if every cell becomes visible from at
least one position visited along the route:

\[
    \forall c\in C_i,\quad
    \exists p\in\pi:\ d_M(p,c)\leq R.
\]

The agent therefore need not visit every spawning cell directly; intermediate
positions traversed while moving between destinations also contribute to
coverage.

\subsection{Greedy Route Construction}

Let $U\subseteq C_i$ contain the cells not yet observed, and let $q$ be the
current planning position. Cells already visible from $q$ are first removed
from $U$. For each reachable candidate destination $v$, an A* path
$\pi(q,v)$ is computed, and its newly covered cells are

\[
    G(q,v)
    =
    \{c\in U \mid \exists p\in\pi(q,v):
    d_M(p,c)\leq R\}.
\]

Candidates are ranked by coverage gained per movement cost,

\[
    u(q,v)=\frac{|G(q,v)|}{|\pi(q,v)|+1},
\]

where $|\pi(q,v)|$ is the number of movement actions. The best path is appended
to the plan, $U$ is updated as $U\leftarrow U\setminus G(q,v)$, and the
planning position becomes $q\leftarrow v$. This process continues until
$U=\emptyset$.

Candidate destinations may be any reachable walkable cells and therefore need
not themselves be spawning cells or belong to the selected cluster.

\subsection{Resumable Search}

A scan can be interrupted when a parcel becomes visible. Instead of discarding
the search progress, the agent stores the scan start time $s_i$, the remaining
uncovered cells, and the latest observation timestamp of each spawning cell.

When the scan resumes, any missing cell $c$ satisfying

\[
    \operatorname{lastObservedAt}(c) \geq s_i
\]

is removed from the remaining set. A new coverage route is then generated from
the agent's current position. Replanning is preferable to following the old
route because parcel collection or delivery may have moved the agent
substantially.

An interrupted cluster remains active, and its completion timestamp is left
unchanged. Only after all remaining cells have been covered is the checkpoint
cleared and $t_i$ set to the current time. The search therefore preserves
progress across interruptions while still guaranteeing complete cluster
coverage.

\subsection{Temporarily Unreachable Clusters}

Giving an active or old cluster priority does not imply that the agent should
remain committed to it when its remaining cells are temporarily unreachable.
Nearby agents and failed movements can introduce temporary walls into the
planning map. If the current cluster were treated as the only admissible
target, one such obstruction could prevent Search from considering every other
cluster, even when another complete coverage route was available.

The agent therefore separates \emph{cluster priority} from
\emph{route feasibility}. Clusters are still considered in the order defined
above---active cluster first, then least recently completed, with distance as
a tie-breaker---but an incomplete route does not immediately terminate the
selection procedure. Instead, the agent preserves that cluster's checkpoint
and evaluates the remaining clusters. The resulting preference order is:

\begin{enumerate}
    \item use the first complete coverage route found in cluster-priority
    order;
    \item if no complete route exists, execute the highest-priority partial
    route that still observes at least some remaining cells;
    \item if every cluster is currently unreachable, keep all checkpoints and
    wait for the next deliberation cycle.
\end{enumerate}

Thus, complete progress on another cluster is preferred to partial progress on
a blocked cluster, while partial progress remains useful as a fallback when no
cluster can currently be completed. Waiting is the final fallback rather than
evidence that Search is permanently infeasible: on the next cycle, sensing is
refreshed, temporary walls are reconsidered, and all clusters are evaluated
again.

This distinction also changes the termination condition. Failure to construct
a route in the presence of temporary walls is classified as a transient block,
not as the absence of every feasible intention. The agent consequently does
not quit merely because the current map snapshot contains no complete route;
it quits only when planning is infeasible without a temporary obstruction that
could disappear on a later cycle.

\subsection{Search Interruption and Persistent Parcel Beliefs}
\label{sec:search-interruption}

Each sensing event reports only the parcels currently visible to the agent.
These observations are merged into a persistent belief state, so the intention
context may also contain remembered parcels outside the current observation
radius. A parcel is therefore not removed merely because it is absent from a
new sensing snapshot.

Previously, Search was interrupted whenever any free parcel existed in the
belief state,

\[
    \operatorname{interrupt}
    \iff
    \operatorname{freeParcelCount} > 0.
\]

This could produce an infinite deliberation loop when a remembered parcel had
already been evaluated as non-viable: Search was selected, but the same parcel
immediately interrupted it before the agent could move.

To prevent this, the agent records the set $K$ of free-parcel identifiers known
when the search plan is created. During execution, Search is interrupted only
when a previously unknown free parcel appears:

\[
    \operatorname{interrupt}
    \iff
    \exists p\in P_{\mathrm{free}}:
    \operatorname{id}(p)\notin K.
\]

Previously evaluated parcels therefore remain in the belief state without
repeatedly cancelling Search, while newly detected parcels still trigger a new
deliberation cycle.\section{Introduction of the Drop-and-Pickup Strategy}

Originally, the agent compared two intentions: \emph{pickup first}, where it
collects a free parcel before delivering all carried parcels, and \emph{drop
only}, where it immediately delivers the parcels it is already carrying.

This comparison used inconsistent planning horizons. The pickup-first score
considered both the carried parcels and the newly collected parcel, whereas
the drop-only score stopped after the first delivery and ignored the possibility
of subsequently collecting a known free parcel.

Let
\[
    V(r,t) = \max(0,r-t)
\]
denote the remaining value of a parcel with current reward $r$ when delivered
after $t$ seconds. Let $C$ be the set of carried parcels, $\mathcal{P}$ the set
of known free parcels, $S$ the current agent position, and $D$ the set of
delivery zones.

\subsection{Pickup-First Intention}

For a free parcel $p \in \mathcal{P}$, the pickup-first alternative considers
the route $S \rightarrow p \rightarrow d$, with $d \in D$. Its estimated
delivery time is

\[
    T_{\mathrm{pickup}}(p)
    = t(S,p) + \min_{d \in D} t(p,d),
\]

giving the utility

\[
    U_{\mathrm{pickup}}(p)
    =
    \sum_{c \in C} V(r_c,T_{\mathrm{pickup}}(p))
    + V(r_p,T_{\mathrm{pickup}}(p)).
\]

Thus, both the currently carried parcels and the newly collected parcel are
evaluated over the same complete pickup-and-delivery horizon.

\subsection{Drop-and-Pickup Intention}

The former drop-only alternative is extended to include the continuation after
the initial delivery. It therefore evaluates routes of the form
$S \rightarrow d_1 \rightarrow p \rightarrow d_2$, where $d_1,d_2 \in D$
and $p \in \mathcal{P}$.

For a candidate first delivery zone $d_1$, define

\[
    T_{\mathrm{drop}}(d_1) = t(S,d_1)
\]

and, for each free parcel $p$,

\[
    T_{\mathrm{cont}}(d_1,p)
    =
    t(S,d_1) + t(d_1,p)
    + \min_{d_2 \in D} t(p,d_2).
\]

The corresponding utility is

\[
    U_{\mathrm{delivery}}(d_1)
    =
    \sum_{c \in C} V(r_c,T_{\mathrm{drop}}(d_1))
    +
    \max_{p \in \mathcal{P}}
    V(r_p,T_{\mathrm{cont}}(d_1,p)).
\]

If no free parcel is available, the continuation term is zero and the
intention reduces naturally to the original drop-only strategy.

The agent can therefore compare two alternatives over comparable planning
horizons:

\[
    S \rightarrow p \rightarrow d
    \qquad\text{vs.}\qquad
    S \rightarrow d_1 \rightarrow p \rightarrow d_2.
\]

This allows the agent to recognize cases in which a delivery zone lies on or
near the route to a valuable parcel. Delivering the carried parcels first may
then preserve their reward while introducing only a small delay before the
subsequent pickup.

Importantly, the complete continuation is used only for scoring. After choosing
a drop-and-pickup intention, the agent executes only the first subgoal,
$S \rightarrow d_1 \rightarrow \text{drop}$, and then deliberates again.
This avoids committing to a parcel that may meanwhile have expired, been
collected by another agent, or become less attractive than a newly available
alternative.

\section{Pathfinding Calculations}

Consider one complete intention-scoring pass with $D=20$ delivery zones
and $P=10$ free parcels. Agent, parcel, and delivery-zone positions are
assumed to be distinct. The agent initially carries one parcel; however,
the number of carried parcels does not affect the number of pathfinding
operations, since their rewards can be evaluated arithmetically once the
delivery time is known.

\subsection{No Cache}

For each pickup-first intention, the agent computes one path from its
current position $S$ to the parcel and one path from that parcel to each
of the $D$ delivery zones. Hence,

\[
C_{\mathrm{pickup}} = P(1+D)
                     = 10(1+20)
                     = 210.
\]

For a delivery-first intention, the agent first computes the path from
$S$ to the selected delivery zone. For each free parcel, it then computes
one path from that delivery zone to the parcel and $D$ paths from the
parcel to the possible final delivery zones. Across all $D$ possible
initial delivery zones,

\[
C_{\mathrm{delivery}}
    = D\bigl[1+P(1+D)\bigr]
    = 20\bigl[1+10(21)\bigr]
    = 4220.
\]

Therefore, without caching,

\[
C_{\mathrm{none}}
    = C_{\mathrm{pickup}} + C_{\mathrm{delivery}}
    = 210 + 4220
    = \boxed{4430}.
\]

\subsection{Directed Cache}

With a directed cache, each unique directed path is computed only once,
so $A\rightarrow B$ and $B\rightarrow A$ require separate entries.
The required paths are $D$ paths from $S$ to the delivery zones, $P$
paths from $S$ to the parcels, $DP$ delivery-to-parcel paths, and $PD$
parcel-to-delivery paths. Thus,

\[
C_{\mathrm{directed}}
    = D + P + 2DP
    = 20 + 10 + 2(20)(10)
    = \boxed{430}.
\]

\subsection{Symmetric Cache}

If the map contains no directional tiles, path lengths are symmetric,

\[
\operatorname{dist}(A,B)=\operatorname{dist}(B,A).
\]

Consequently, delivery-to-parcel and parcel-to-delivery paths can share
the same cache entry. Only $D$ agent-to-delivery, $P$ agent-to-parcel,
and $DP$ parcel--delivery paths are required:

\[
C_{\mathrm{symmetric}}
    = D + P + DP
    = 20 + 10 + (20)(10)
    = \boxed{230}.
\]

\subsection{Comparison}

\begin{table}[t]
    \centering
    \begin{tabular}{lr}
        \toprule
        \textbf{Strategy} & \textbf{A* calculations} \\
        \midrule
        No cache        & 4430 \\
        Directed cache  & 430  \\
        Symmetric cache & 230  \\
        \bottomrule
    \end{tabular}
    \caption{Number of A* computations for $D=20$ delivery zones and
    $P=10$ free parcels.}
    \label{tab:pathfinding}
\end{table}

Relative to the uncached implementation, the directed cache reduces the
number of A* computations by a factor of

\[
\frac{4430}{430} \approx 10.3,
\]

while the symmetric cache achieves a reduction factor of

\[
\frac{4430}{230} \approx 19.3.
\]

The symmetric cache therefore requires approximately $47\%$ fewer
pathfinding computations than the directed cache ($230$ versus $430$),
provided that path costs are symmetric.

\section{Collision Avoidance under Asynchronous Sensing}
\label{sec:collision-avoidance}

Deliveroo.js is a real-time rather than turn-based environment. Movement
requests may reach the server at any instant, while sensing information is
generated only at discrete frames. Consequently, an observation may already
be stale when the next movement request is processed.

\subsection{Movement and Tile Locking}

When an agent moves from a source tile $x$ to an adjacent destination $y$, the
source is initially locked by occupancy. The server validates and locks $y$,
updates the agent to an intermediate fractional position, waits for the
movement duration, and finally places it at $y$ while releasing $x$.

For example, during a movement from $4$ to $5$, the observed position evolves
as $4 \rightarrow 4.6 \rightarrow 5$. At the intermediate position both tiles
$4$ and $5$ are locked, although the spatial registry associates the agent
with the destination cell rather than with two cells simultaneously.

A fractional coordinate $p$ therefore identifies two temporarily unavailable
cells:

\[
B(p)=
\begin{cases}
\{\lfloor p\rfloor,\lceil p\rceil\}, & p \text{ fractional},\\
\{p\}, & p \text{ integer}.
\end{cases}
\]

Since movement occurs along one axis at a time, the same reasoning applies to
the corresponding coordinate in two dimensions.

\subsection{Frames and Stale Sensing}

A server frame is an observation and synchronization boundary, not an action
turn. Tile locks and movement requests may occur between two consecutive
frames, whereas a sensing snapshot is generated only when the corresponding
sensor is updated.

Thus, a tile reported as free at frame $k$ may be locked by another agent
before our movement request reaches the server. The request can consequently
fail even though the preceding observation was correct when generated. A
later frame $k+1$ may then reveal the movement that caused the conflict.

The sensing state should therefore be interpreted as recent information rather
than as a guarantee of future accessibility. The server acknowledgement of a
movement request remains the authoritative indication of whether the action
succeeded.

\subsection{Absence of a Movement Cooldown}

There is no guaranteed idle interval between consecutive movements. Once a
movement completes, the server places the agent at its destination, releases
the source tile and action mutex, and sends the acknowledgement. A subsequent
movement may begin as soon as the next request arrives.

Network latency introduces a small practical delay, but its duration is
variable and cannot be used as a reliable synchronization mechanism. In
particular, waiting for one nominal movement duration after observing an
intermediate position is insufficient: the observation itself is delayed, and
the other agent may immediately start another movement after completing the
current one. Collision avoidance must therefore rely on observed state rather
than fixed timers.

\subsection{Conservative Collision-Avoidance Policy}

Fractional positions provide information about movement direction. Observing
an agent at $4.6$, for example, indicates a committed transition from $4$ to
$5$; until this movement completes, both cells are unavailable and the agent
cannot reverse direction. In contrast, observing it at the integer position
$5$ establishes only its current occupancy: it may stop, continue toward $6$,
or immediately reverse toward $4$.

The adopted policy therefore:

\begin{enumerate}
    \item maintains the latest sensed position of each agent;
    \item blocks both cells associated with a fractional position;
    \item treats an integer-positioned agent's current cell as occupied and
          neighbouring cells as potentially unsafe;
    \item checks these dynamic obstacles before each planned movement;
    \item suspends and replans when another agent enters the relevant safety
          region;
    \item enters a recently vacated cell only after sensing confirms that the
          other agent is committed to moving farther away; and
    \item invalidates the remaining plan immediately if a movement request
          returns \texttt{false}.
\end{enumerate}

For instance, if the desired tile is $4$, an agent observed at $4.6$ is still
locking both $4$ and $5$. At position $5$, tile $4$ may already be free, but an
immediate reversal remains possible. Once the agent is observed at $5.6$, it
is committed to moving toward $6$, making tile $4$ substantially safer to
enter. The path is then recomputed using the most recent sensing state.

This policy cannot eliminate all races because sensing and acting are not
atomic. It nevertheless reduces collision penalties by exploiting fractional
positions as evidence of committed movement and by avoiding synchronization
based solely on timing estimates.

Passive waiting may itself cause deadlock when two controlled agents adopt the
same strategy. In that case, a deterministic right-of-way rule, for example
one based on agent identifiers, can break the symmetry. Otherwise, waiting
should be bounded and followed by another attempt to find an alternative
route.

\section{Mechanisms for Handling Other Agents}
\label{sec:agent-handling}

The agent uses a conservative collision-avoidance strategy. Before each
movement, the intended destination is checked against the latest observations
of nearby agents. When a conflict is detected, the agent either waits for
additional sensing information or temporarily excludes the contested cell and
replans.

\subsection{Movement Reconstruction and Guarding}

During movement, the server may report an agent at an intermediate fractional
coordinate, approximately $0.6$ of the way between two adjacent cells. This
allows the source and destination of the movement to be reconstructed and,
therefore, to distinguish whether the other agent is moving toward the intended
cell, away from it, or along an unrelated trajectory.

Immediately before issuing a movement command, the destination is checked
against the latest sensing state. If another agent is moving toward the
destination, the agent waits until it has arrived, left the cell, and revealed
its subsequent direction. If the other agent is instead moving away from the
destination, the guard waits until that departure is confirmed and ensures
that it does not immediately reverse. Unrelated movements do not delay the
current plan.

Waiting is driven by new sensing snapshots rather than repeated deliberation.
Once the destination is considered safe, execution of the original movement
continues.

\subsection{Incomplete Observations}

Intermediate movement states are not guaranteed to appear in every sensing
snapshot. The guard therefore does not require the complete sequence of
fractional and integer positions to be observed.

If a later observation places the tracked agent farther from the contested
cell or moving along an unrelated trajectory, its departure is considered
confirmed. If the agent disappears from sensing, the guard checks the area
covered by the latest authoritative snapshot. When the destination lies inside
the observed area but the other agent is absent, the cell is considered empty;
otherwise, the guard continues waiting.

\subsection{Stationary Agents and Oscillation}

A nearby agent that remains stationary in a relevant position is tolerated for
at most

\[
    T_{\mathrm{stationary}} = 2T_{\mathrm{move}}.
\]

If it remains stationary beyond this interval, its occupied cell is converted
into a temporary wall and deliberation is restarted. The same rule applies
when an agent leaves the contested cell but then stops in a position that
prevents the guard from safely confirming continuation.

Similarly, if an agent moves away from a contested cell and then returns, the
behaviour is classified as oscillation. Rather than waiting indefinitely for a
repeated back-and-forth movement to terminate, the contested cell is
temporarily blocked and a new plan is generated.

\subsection{Temporary Walls and Replanning}

Temporary walls are considered both during intention evaluation and during
path construction. Consequently, pickup, delivery, and exploration plans are
all generated without crossing the blocked cell.

A temporary wall remains active for one complete deliberation cycle and is
then removed. This prevents transient conflicts from permanently excluding
cells that may later become usable, particularly when no alternative route
exists.

\subsection{Failed-Movement Fallback}

A race condition remains possible when two agents observe the same free cell
and submit movement requests before receiving another sensing update. The
server resolves the conflict by allowing the first request to lock the cell
and rejecting the other.

If a movement request fails, the agent keeps its authoritative position
unchanged, leaves the failed action unexecuted, temporarily blocks the attempted
destination, invalidates the remaining plan, and immediately deliberates again.

This fallback cannot prevent the penalty associated with losing the race, but
it prevents the agent from continuing under the false assumption that the
movement succeeded.

\subsection{Diagnostic Logging}

The behaviour is exposed through structured diagnostic messages:
\texttt{AGENT ENCOUNTER} identifies an initial conflict,
\texttt{AGENT MOVEMENT} reports relevant trajectory changes,
\texttt{AGENT CLEAR} indicates that execution may continue,
\texttt{AGENT REPLAN} records stationary or oscillating conflicts, and
\texttt{MOVE FAILED} reports server-rejected movements.

Each deliberation also logs any temporary wall used during intention evaluation
and path planning.

\newpage

\section{Unified Pickup-and-Delivery Planning}
\label{sec:unified-planning}

The previous intention-selection mechanism evaluates pickup and delivery using
different and relatively short planning horizons. A pickup intention considers
collecting one free parcel and subsequently delivering it, whereas a delivery
intention considers delivering the currently carried parcels followed by at
most one pickup.

This can lead to locally optimal but globally poor decisions. For example, an
agent may select a nearby high-value parcel even though collecting another
parcel first would allow several parcels to be delivered together with a
greater total reward. More generally, routes such as

\[
    p_1 \rightarrow p_2 \rightarrow d
    \qquad\text{and}\qquad
    d \rightarrow p_1 \rightarrow p_2 \rightarrow d'
\]

cannot be compared fairly when the two initial actions are evaluated with
different amounts of look-ahead.

\subsection{Unified Planning Model}

Pickup and delivery are therefore treated as transitions of the same routing
problem rather than as independent intention types. A planning state contains
the agent position, the sets of free and carried parcels, the elapsed time, and
the accumulated reward.

From each state, two types of transition are possible:

\begin{enumerate}
    \item move to a known free parcel and collect it;
    \item move to a delivery cell and deliver all currently carried parcels.
\end{enumerate}

The resulting search can represent arbitrary combinations such as
$p_1\rightarrow p_2\rightarrow d$, $d\rightarrow p_1\rightarrow d'$ or
$p_1\rightarrow d\rightarrow p_2\rightarrow p_3\rightarrow d'$ within the
chosen planning horizon. Parcel rewards are evaluated at their estimated
delivery times, so delaying a delivery is reflected directly in the route
score.

Pickup and delivery consequently share the same route evaluator and search
budget. The first transition of the highest-scoring route determines the next
intention: a parcel transition produces a pickup intention, while a delivery
transition produces a delivery intention.

Only this first route segment is executed. The agent then deliberates again,
since parcels may appear, decay, disappear, or be collected by another agent
while the planned continuation is being executed.

This problem resembles the travelling-salesperson problem, but differs from
the classical formulation because parcel visits are optional, delivery cells
may be revisited, multiple parcels can be delivered simultaneously, rewards
decay over time, and carrying capacity may constrain the route. It is therefore
better viewed as a time-dependent pickup-and-delivery routing problem.

\subsection{Delivery-Cell Candidate Reduction}

Considering every delivery cell at every stage can greatly increase the search
space. Selecting only the delivery cell closest to each parcel is attractive
but unsafe, since the best intermediate delivery point depends on both the
previous and subsequent locations.

For locations $u$ and $v$, the most useful intermediate delivery cell can
instead be defined as

\[
    d^*(u,v)
    =
    \arg\min_{d\in D}
    \bigl[d(u,d)+d(d,v)\bigr].
\]

A delivery cell that is closest to neither endpoint may still minimize this
combined cost by lying directly along the route between them.

A reduced candidate set can therefore include the delivery cell closest to
the current agent position, the closest terminal delivery cell for each
relevant parcel, and $d^*(u,v)$ for relevant pairs of locations. Since several
pairs often select the same cell, duplicate candidates are removed before
route search.

\subsection{Search Complexity}

The planning horizon can limit the number of pickups without discarding
parcels in advance. For example, with seven available parcels and at most five
pickups per candidate route, the number of possible ordered parcel sequences is

\[
    \sum_{k=1}^{5} P(7,k)
    = 7 + 42 + 210 + 840 + 2520
    = 3619.
\]

This remains practical if shortest-path distances are computed once between
the agent position, parcel locations, and candidate delivery cells. Candidate
routes can then be evaluated using cached distances rather than invoking A*
for every transition of every route.

For small parcel sets, an exact branch-and-bound or dynamic-programming search
can be used. Branch-and-bound discards partial routes whose best possible
continuation cannot improve the current solution. When the number of parcels
becomes larger, beam search provides a bounded alternative by retaining only a
fixed number of the most promising partial routes at each search depth.

Constraint-based formulations such as CP-SAT or SMT are also possible, but
explicit route search is more natural for the relatively small and dynamic
planning problems considered here.

\subsection{Proposed Implementation}

The unified planner therefore operates on a compact graph whose nodes represent
the current agent position, known parcels, and selected delivery cells.
Shortest-path distances between these nodes are cached, while pickup and
delivery transitions are explored within a common search procedure.

Each candidate route is scored using the expected parcel rewards at their
respective delivery times. The planner returns the highest-scoring route, but
only its first transition is committed to execution. Subsequent sensing then
triggers a new deliberation, combining longer-horizon decision making with the
reactivity required by the dynamic environment.

\section{PDDL-Based Planning with Fast Downward}
\label{sec:pddl-fast-downward}

The agent uses the Planning Domain Definition Language (PDDL) for navigation
problems that require reasoning about changes to the environment rather than
ordinary pathfinding alone. PDDL describes the current state, the available
actions, their preconditions and effects, and the desired goal state. The
resulting planning problem is solved locally using Fast Downward.

PDDL should not be confused with DPLL: PDDL is a language for representing
automated planning problems, whereas DPLL is an algorithm for propositional
satisfiability.

\subsection{Role of PDDL}

During deliberation, the agent first attempts to construct a plan using its
ordinary pathfinding mechanism. This approach is inexpensive and is preferred
whenever the selected target can be reached without changing the environment.

If no ordinary action sequence can be constructed, the agent invokes the PDDL
planner. The primary use case is navigation obstructed by movable crates. In
this setting, shortest-path search alone is insufficient because pushing a
crate changes the map and may enable or disable subsequent movements.

PDDL therefore models actions with explicit state-changing effects. A crate
push, for example, requires the agent to be adjacent to the crate and the cell
behind it to be free; its effects update both the agent and crate positions.

In the current implementation, PDDL is used as a fallback when the ordinary
action builder returns an empty plan. Since an empty plan may also indicate
that the goal is already satisfied, intentions check their goal condition
before requesting PDDL planning whenever possible.

\subsection{Domain and Problem Representation}

The planning model is separated into a static domain and a dynamically generated
problem.

The domain defines the object types \texttt{position}, \texttt{agent},
\texttt{parcel}, and \texttt{crate}, together with predicates representing
agent, parcel, and crate positions, map adjacency, pickup and delivery cells,
carried parcels, and completed deliveries. It also defines movement, crate
pushing, pickup, and delivery actions through their preconditions and effects.

For each planning request, the problem is generated from the agent's current
belief state. It contains the valid map cells and adjacency relations, the
current agent position, known parcel and crate positions, and a goal determined
by the selected intention. A search intention requires reaching a target cell;
a pickup intention requires reaching and carrying a parcel; and a delivery
intention requires delivering the currently carried parcels.

This separation allows the same domain model to be reused while only the
problem state and goal change between deliberation cycles.

\subsection{Fast Downward Integration}

Fast Downward is executed locally as a child process. For each planning
request, the agent generates the PDDL problem, writes the domain and problem to
an isolated temporary directory, and invokes Fast Downward using the
\texttt{lama-first} configuration. Fast Downward translates the task into its
internal SAS+ representation, searches for a plan, and writes the resulting
action sequence to a plan file.

Conceptually, the invocation is:

\begin{verbatim}
fast-downward.py --alias lama-first \
    --overall-time-limit 30s \
    --plan-file sas_plan \
    domain.pddl problem.pddl
\end{verbatim}

The \texttt{lama-first} configuration is satisficing: it aims to return a valid
plan quickly rather than prove optimality. This is appropriate in an online
environment, where the world may change before a computationally expensive
optimal plan can be executed.

Planning is limited to 30 seconds. If the planner determines that no plan
exists, the result is converted into an empty action sequence. Timeouts,
malformed PDDL, and process failures are instead treated as planning errors.

\subsection{Plan Conversion and Execution}

The Fast Downward output consists of grounded actions, for example

\begin{verbatim}
(move-right ag_29be11 t_5_4 t_5_5)
(pick-up ag_29be11 parcel_p1170 t_5_5)
\end{verbatim}

The parser extracts each action and its arguments and converts them into the
corresponding Deliveroo actions. Movement and crate-pushing actions map directly
to movement commands, while pickup and delivery actions additionally require
removing the prefixes introduced by the PDDL object representation.

The converted actions are inserted into the current plan and executed
sequentially. If an action fails or new sensing information invalidates the
plan, execution is interrupted and a new deliberation cycle begins.

\subsection{Local Planning versus a Remote API}

An earlier implementation submitted PDDL problems to a remote planning service
through HTTP. Although this removed the need for a local planner installation,
it introduced network latency, server queueing, external availability
dependencies, and less predictable timeouts.

Fast Downward provides a native command-line interface, making local execution
a more direct integration. It eliminates network and queueing delays, fixes the
planner version used by the project, allows direct control over time limits and
output files, and improves reproducibility because the same domain and problem
files can be executed independently.

The local approach requires Fast Downward to be installed and consumes local
computational resources. Nevertheless, its lower and more predictable latency,
reproducibility, and independence from an external service make it better
suited to an online agent that may need to replan repeatedly as the environment
changes.

\end{document}
